\documentclass[11pt]{article}

\usepackage[margin=1in]{geometry}
\usepackage{amsmath,amssymb,bm,mathtools}
\usepackage{booktabs}
\usepackage[numbers,sort&compress]{natbib}
\usepackage[colorlinks=true,allcolors=blue]{hyperref}
\usepackage[nameinlink,noabbrev]{cleveref}
\newcommand{\F}{\mathbf{F}}
\newcommand{\Fb}{\bar{\mathbf{F}}}
\newcommand{\A}{\mathbf{A}}
\newcommand{\I}{\mathbf{I}}
\newcommand{\C}{\mathbb C}
\newcommand{\D}{\mathbb D}
\newcommand{\Hh}{\mathbb H}
\newcommand{\Pvol}{\mathbb P_{\mathrm{vol}}}
\newcommand{\Pdev}{\mathbb P_{\mathrm{dev}}}
\newcommand{\tr}{\operatorname{tr}}
\newcommand{\sym}{\operatorname{sym}}
\newcommand{\dd}{\,\mathrm d}
\newcommand{\pder}[2]{\frac{\partial #1}{\partial #2}}
\newcommand{\atp}{\big|_{p}}
\newcommand{\ate}{\big|_{\bm e}}
\providecommand{\doi}[1]{\href{https://dx.doi.org/#1}{#1}}

\title{An anisotropic Biot tensor from finite-deformation mixture kinematics}
\author{John T. Foster}
\date{}
\hypersetup{pdftitle={An anisotropic Biot tensor from finite-deformation mixture kinematics},pdfauthor={John T. Foster}}

\begin{document}
\maketitle

\begin{abstract}
Pore pressure can couple directionally to stress in an anisotropic porous
solid, yet tensorial Biot coefficients are commonly introduced as incremental
quantities.  We derive a finite-deformation Biot tensor for a saturated,
one-solid/one-fluid mixture as the fixed-pressure sensitivity of true mineral
volume to skeleton deformation.  The derivation retains the mineral volume
ratio as an equilibrium variable, so the resulting tensor is a
pressure- and deformation-dependent state function rather than only a
reference-state coefficient.  Under isotropic distention, mixture kinematics
fix the true-mineral shape from the skeleton deformation and mineral volume.
A generalized-Hooke reduction then identifies the mixed
mineral-volume--skeleton-strain tangent that controls the small-strain tensor.
It shows precisely how mineral anisotropy enters: only the
volumetric--deviatoric block \(\Pdev:\Hh:\I\) contributes a directional
mineral-stiffness term.  Mineral anisotropy confined to the deviatoric block
does not.  An explicit volumetric calibration recovers
\(1-K/K_s\).  The analysis also establishes that a drained skeleton tangent
and a mineral-volume tangent do not determine the pressure coupling without
the mixed tangent.
\end{abstract}

\noindent\textbf{Keywords:} finite-deformation poromechanics; anisotropic
elasticity; Biot tensor; mineral stiffness; Legendre transform

\section{Introduction}
\label{sec:introduction}

The linear theory of a fluid-saturated anisotropic porous solid was formulated
by Biot \citep{biot1955anisotropic}.  Its pressure--stress coupling requires,
in addition to the fourth-order elastic tensor, a symmetric second-order
effective-stress (Biot--Willis) tensor and a storage coefficient
\citep{biotwillis1957,thompsonwillis1991}.  This distinction is consequential:
anisotropic elastic stiffness alone does not determine the directional
pressure coupling.  Micromechanical interpretations and practical
identification procedures have made this point explicit under assumptions
about the solid constituents and pore structure \citep{cheng1997}, while
crack-based analyses show how pore fabric can generate anisotropic Biot--Willis
and Skempton tensors \citep{wong2017}.

The experimental literature also motivates retaining tensorial coupling for
layered rocks.  Dynamic measurements on VTI shale establish direction- and
stress-dependent elastic stiffnesses \citep{sviridov2017}, and a more complete
set of drained, pore-pressure, and deviatoric tests on Callovo--Oxfordian
claystone found a strongly anisotropic strain response and a markedly
anisotropic Skempton response, even though the back-calculated Biot-tensor
anisotropy was small in that material \citep{braun2020}.  At the fracture
scale, hydraulic fracturing in anisotropic shale is a potential application
for poromechanical models with tensorial pressure coupling
\citep{chaubazantsu2016}.

The Biot coefficient measures the pressure contribution to total stress.  In
anisotropic theories it is natural to represent that contribution by a tensor,
but the tensor is often specified as an incremental closure.  For example,
Zhao and Borja \citep{zhaoborja2020} obtain an anisotropic small-strain Biot
tensor from the drained elastic tangent and a scalar intrinsic bulk modulus.
Such a construction captures directional incremental response, but it does
not itself prescribe the mineral-volume equilibrium needed to follow the
coefficient through finite deformation and finite pressure.

In our earlier work \citep{fosterxu2025}, mixture kinematics,
intrinsic-density conjugacy, and a pressure Legendre transform gave the
fixed-pressure sensitivity of the intrinsic solid specific volume,
\(\bar v_s\), to skeleton deformation.  Its scalar specialization assumes
\(\bar v_s=\bar v_s(J)\), and hence \(\phi_s=\phi_s(J)\).  Once the solid
volume fraction depends on the full deformation gradient,
\(\phi_s=\phi_s(\F,\ldots)\), the true mineral volume depends on more than
the skeleton Jacobian and its fixed-pressure sensitivity is tensorial.

We develop that finite-pressure tensor for one deformable solid and one fluid
at local pressure equilibrium.  The general result follows from the
distention--true-deformation kinematics and the Legendre transform, without
selecting an elastic law.  It therefore identifies a state-dependent tensor
\(\bm B(\F,p)\), not merely a linearized coefficient.  We then specialize to
isotropic distention and a generalized-Hooke energy.  This reduction separates
the skeleton stiffness, mineral stiffness, and the mixed tangent that relates
skeleton strain to mineral volume.  The mixed tangent is indispensable: the
two individual stiffnesses do not determine how mineral volume changes at
fixed pore pressure.

The special cases make the mineral-anisotropy mechanism explicit.  A mineral
stiffness contributes directly to the directional Biot tensor only through
its volumetric--deviatoric coupling; anisotropy confined to the deviatoric
block leaves that direct contribution isotropic.  A volumetric calibration
then recovers the classical coefficient \(1-K/K_s\)
\citep{biotwillis1957}.  These reductions connect the finite-pressure
construction to the familiar scalar limit while isolating the constitutive
information required for directional coupling.

\section{Kinematics and thermodynamic setting}
\label{sec:kinematics}

Let \(\F\) map the solid skeleton reference configuration to its current
configuration, and let \(J=\det\F>0\) be the corresponding skeleton volume
ratio.  The solid partial density is
\(\rho_s=\phi_s\bar\rho_s\), where \(\phi_s\) is the solid volume fraction
and \(\bar\rho_s\) is the intrinsic density of the mineral solid.  For a
uniform reference state, solid mass conservation gives
\begin{equation}
 J\rho_s=\rho_{s0}=\phi_{s0}\bar\rho_{s0}.
 \label{eq:solid-mass}
\end{equation}
The true mineral volume ratio is
\begin{equation}
 \bar J\equiv\frac{\bar\rho_{s0}}{\bar\rho_s},
 \qquad
 \phi_s=\frac{\phi_{s0}\bar J}{J},
 \qquad \phi_f=1-\phi_s.
 \label{eq:distention}
\end{equation}
Thus \(\bar J\) is the mineral volume ratio, \(\phi_f\) is the fluid volume
fraction, and \(J/\bar J\) is the distention ratio of the skeleton.  In other
words, \(J\) measures the change of the porous skeleton volume, whereas
\(\bar J\) measures the change of volume of the mineral itself.  We separate
these two changes using the
distention--true-deformation construction emphasized by Drumheller
\citep{drumheller2000},
\begin{equation}
 \F=\A\Fb,
 \qquad \det\A=\frac{J}{\bar J},
 \qquad \det\Fb=\bar J.
\label{eq:true-mineral-jacobian}
\end{equation}
Here \(\A\) describes distention and \(\Fb\) describes true mineral
deformation.  The determinant relation in
\eqref{eq:true-mineral-jacobian} is enough to track volume.  We use the
isotropic-distention specialization below, which also fixes the shape part of
the true mineral deformation once \(\F\) and \(\bar J\) are known.
For the nonreacting, one-solid/one-fluid specialization considered here, an
isotropic distention \(\A=a\I\), with
\(a=(J/\bar J)^{1/3}\), makes \(\F=a\Fb\) kinematically equivalent, up to
the naming and ordering of factors, to the isotropic swelling--elastic
multiplicative split used in Flory-type network theories \citep{flory1961}.
This is a kinematic equivalence only: it neither transfers Flory's mixing free
energy nor supplies the present pressure-conjugacy or constitutive assumptions.

Let \(q=\ln\bar J\) be the mineral volume coordinate.  At the fixed
temperature considered here, let \(\psi_s(\F,q)\) be the solid specific
Helmholtz free energy and define the mixture-reference free-energy density by
\(W(\F,q)=\rho_{s0}\psi_s(\F,q)\).  The normalization used
in our earlier work \citep{fosterxu2025} takes the intrinsic solid specific
volume as \(\bar v_s=1/\bar\rho_s=\bar v_{s0}\bar J\), and
\(\rho_{s0}\bar v_{s0}=\phi_{s0}\).  Following Foster and Xu
\citep{fosterxu2025}, \(p\) denotes the fluid pore pressure.  The isothermal pressure conjugacy
\(\partial\psi_s/\partial\bar v_s=-p\) then gives the following
mixture-reference relation.  Structural tensors and history variables are
suppressed and held fixed below.  Local pressure equilibrium determines the
mineral volume coordinate:
\begin{equation}
 \left.\pder{W}{q}\right|_{\F}=-\phi_{s0}p\bar J.
 \label{eq:mineral-coordinate-equilibrium}
\end{equation}
Equivalently, the volume derivative at fixed \(\F\) is
\begin{equation}
 \left.\pder{W}{\bar J}\right|_{\F}=-\phi_{s0}p.
 \label{eq:pressure-conjugacy}
\end{equation}
Consider an already-deformed specimen that is rotated rigidly in space, with
its shape and pore pressure left unchanged.  This change of laboratory
orientation cannot affect the stored energy or the mineral volume selected by
local equilibrium.  Thus, \(\bar J\) does not depend on how the specimen is
oriented.

\section{Finite-deformation Biot tensor}
\label{sec:finite-biot}

Assume that \eqref{eq:mineral-coordinate-equilibrium} can be solved locally
for \(q=q(\F,p)\); specifically, assume that the derivative of
\(W_q+\phi_{s0}p\bar J\) with respect to \(q\) is nonzero.  Define the
reduced fixed-pressure energy and its
Legendre transform by
\begin{subequations}
\label{eq:fixed-pressure-energies}
\begin{align}
 W^{\prime\prime}(\F,p)&=W\bigl(\F,q(\F,p)\bigr),
 \label{eq:reduced-energy}\\
 W^{\prime}(\F,p)&=W^{\prime\prime}(\F,p)+\phi_{s0}p\,
 \bar J\bigl(q(\F,p)\bigr).
 \label{eq:legendre-energy}
\end{align}
\end{subequations}
In our earlier work \citep[eq.~(33)]{fosterxu2025}, we transformed the specific
internal energy according to
\(\widetilde e^s(\F,p)=e^s(\F,\bar v_s)+p\bar v_s\).  Equation
\eqref{eq:legendre-energy} is its isothermal Helmholtz counterpart.  Replacing
\(e^s\) by \(\psi_s\), multiplying by \(\rho_{s0}\), and using
\(\bar v_s=\bar v_{s0}\bar J\) together with
\(\rho_{s0}\bar v_{s0}=\phi_{s0}\) gives
\(W+\phi_{s0}p\bar J=\rho_{s0}(\psi_s+p\bar v_s)\).  Evaluating this
identity at the local solution \(q(\F,p)\) gives
\eqref{eq:legendre-energy} per unit reference mixture volume.
The first Piola stresses conjugate to \(\F\), respectively at a fixed mineral
state and at a fixed pore pressure, are
\begin{subequations}
\label{eq:effective-stresses}
\begin{align}
 \bm P^{\prime}&=\left.\pder{W}{\F}\right|_{q}
 =\left.\pder{W^{\prime}}{\F}\right|_p,
 \label{eq:fixed-mineral-state-stress}\\
 \bm P^{\prime\prime}&=\left.\pder{W^{\prime\prime}}{\F}\right|_p.
 \label{eq:fixed-pressure-stress}
\end{align}
\end{subequations}
The second equality in the first definition follows because the change in
\(q\) cancels through
\eqref{eq:mineral-coordinate-equilibrium}.  Differentiating
\eqref{eq:legendre-energy} gives
\begin{equation}
 \bm P^{\prime}=
 \bm P^{\prime\prime}+
 \phi_{s0}p\left.\pder{\bar J}{\F}\right|_p.
 \label{eq:piola-transform-general}
\end{equation}

The total mixture first Piola stress is
\(\bm P=\bm P^{\prime}-pJ\F^{-T}\).  The first term contains the stress
transmitted through the skeleton and mineral state, while the second is the
usual pore-pressure traction.  We define the finite-deformation Biot tensor in
the current configuration by
\begin{equation}
 \bm B
 =\I-\frac{\phi_{s0}}{J}
 \left.\pder{\bar J}{\F}\right|_p\F^T.
 \label{eq:finite-biot-tensor}
\end{equation}
Equations \eqref{eq:piola-transform-general} and \eqref{eq:finite-biot-tensor}
then give the familiar effective-stress form
\begin{equation}
 \bm\sigma=\bm\sigma^{\prime\prime}-p\bm B,
 \label{eq:total-stress-transform}
\end{equation}
where \(\bm\sigma^{\prime\prime}=J^{-1}\bm P^{\prime\prime}\F^T\).
Under the rotation-invariance assumption stated above, a superposed rigid
rotation \(\bm Q\) gives \(\bm B(\bm Q\F,p)=\bm Q\bm B(\F,p)\bm Q^T\),
and the same relation holds for \(\bm\sigma^{\prime\prime}\).  For comparison,
suppose that the mineral volume responds only to skeleton volume and not to
skeleton shape.  Then \(\bar J\) depends on \(\F\) only through \(J\), and
\begin{equation}
 \bm B=\left(1-\phi_{s0}\left.\pder{\bar J}{J}\right|_p\right)\I,
 \label{eq:scalar-recovery}
\end{equation}
which is the scalar finite-deformation result derived in our earlier work
\citep[eq.~(39)]{fosterxu2025}.  The normalization used there gives the same
expression exactly: the reference specific volume is \(v_{s0}=1/\rho_{s0}\), where
\(\rho_{s0}=\phi_{s0}\bar\rho_{s0}\), whereas the intrinsic specific volume
is \(\bar v_s=\bar J/\bar\rho_{s0}\).  Hence
\begin{equation}
 \frac{\bar v_s}{v_{s0}}=\phi_{s0}\bar J,
 \qquad
 B=1-\phi_{s0}\left.\pder{\bar J}{J}\right|_p.
 \label{eq:foster-xu-normalization-bridge}
\end{equation}
This is equation~(39) of \citet{fosterxu2025}, expressed using the present
volume ratio.

The tensor in \eqref{eq:finite-biot-tensor} is symmetric under the stated
rotation-invariance assumption.  To show this, take a superposed infinitesimal rotation
\(\bm Q=\I+\bm\Omega\), \(\bm\Omega^T=-\bm\Omega\).  Invariance of
\(\bar J\) gives
\begin{equation}
 \left.\pder{\bar J}{\F}\right|_p:(\bm\Omega\F)=0,
 \qquad\text{so}\qquad
 \left.\pder{\bar J}{\F}\right|_p\F^T
 \ \text{is symmetric}.
 \label{eq:biot-tensor-symmetry}
\end{equation}

\section{Finite-deformation Hencky realization}
\label{sec:finite-hencky}

The kinematic result above does not require a particular elastic law.  We now
give one objective finite-deformation realization that respects the
isotropic-distention split.  Let \(\F=\bm R\bm U\) and
\(\Fb=\bar{\mathbf R}\bar{\mathbf U}\) be right polar decompositions, and
introduce the material logarithmic strains
\begin{equation}
 \bm\varepsilon=\log\bm U,
 \qquad
 \bm\eta=\log\bar{\mathbf U}.
 \label{eq:hencky-strain-coordinates}
\end{equation}
Because \(\A=a\I\), \(\bar{\mathbf U}=a^{-1}\bm U\).  Thus \(\bm\eta\)
is not an independent state variable: it is fixed by \(\bm\varepsilon\) and
\(q\).  Introduce the volumetric and deviatoric projectors on symmetric
tensors,
\begin{equation}
 \Pvol=\frac13\I\otimes\I,
 \qquad
 \Pdev=\mathbb I_{\mathrm{sym}}-\Pvol.
 \label{eq:volumetric-deviatoric-projectors}
\end{equation}
Then
\begin{equation}
 \bm\eta=\Pdev:\bm\varepsilon+\frac{q}{3}\I.
 \label{eq:constrained-mineral-hencky-strain}
\end{equation}
In particular, \(\tr\bm\eta=q=\ln\bar J\) exactly.  Let
\(W_{\mathrm H}(\bm\varepsilon,q)\) be any objective Hencky energy per unit
reference mixture volume, with structural tensors understood in the material
frame.  At prescribed pore pressure, the scalar mineral-volume equilibrium is
\begin{equation}
 \frac{\partial W_{\mathrm H}}{\partial q}
 +\phi_{s0}p\exp(q)=0.
 \label{eq:finite-hencky-mineral-equilibrium}
\end{equation}
Define the scalar fixed-pressure mineral tangent by
\begin{equation}
 h_p=\frac{\partial^2 W_{\mathrm H}}{\partial q^2}
 +\phi_{s0}p\exp(q).
 \label{eq:finite-hencky-mineral-tangent}
\end{equation}
For \(h_p\ne0\), differentiation of
\eqref{eq:finite-hencky-mineral-equilibrium} at fixed pressure gives
\begin{equation}
 \left.\frac{\partial q}{\partial\bm\varepsilon}\right|_p
 =-h_p^{-1}\frac{\partial^2 W_{\mathrm H}}
 {\partial q\,\partial\bm\varepsilon}.
 \label{eq:finite-hencky-mineral-sensitivity}
\end{equation}
Since \(\dd\bar J=\bar J\,\dd q\),
\eqref{eq:finite-hencky-mineral-sensitivity} and the chain rule from
\(\bm\varepsilon(\F)=\log\bm U\) give
\begin{equation}
 \left.\pder{\bar J}{\F}\right|_p
 =-\bar J h_p^{-1}
 \frac{\partial^2 W_{\mathrm H}}
 {\partial q\,\partial\bm\varepsilon}:
 \frac{\partial\bm\varepsilon}{\partial\F}.
 \label{eq:finite-hencky-volume-sensitivity}
\end{equation}
Here \(\partial\bm\varepsilon/\partial\F\) is the fourth-order tangent
that maps \(\dd\F\) to \(\dd\bm\varepsilon\).  Substitution in
\eqref{eq:finite-biot-tensor} yields
\begin{equation}
 \bm B=\I+\frac{\phi_{s0}\bar J}{J}h_p^{-1}
 \left[\frac{\partial^2 W_{\mathrm H}}
 {\partial q\,\partial\bm\varepsilon}:
 \frac{\partial\bm\varepsilon}{\partial\F}\right]\F^T.
 \label{eq:finite-hencky-biot-tensor}
\end{equation}
At the zero-pressure reference state,
\(\bm\varepsilon=\bm\eta=\bm0\), \(J=\bar J=1\),
\(\left.\partial\bm\varepsilon/\partial\F\right|_{\F=\I}
=\mathbb I_{\mathrm{sym}}\), the identity map on symmetric tensors.
Equation~\eqref{eq:finite-hencky-biot-tensor} therefore becomes
\begin{equation}
 \bm B_0=\I+\phi_{s0}h_0^{-1}
 \left.\frac{\partial^2 W_{\mathrm H}}
 {\partial q\,\partial\bm\varepsilon}\right|_0,
 \qquad
 h_0=\left.\frac{\partial^2 W_{\mathrm H}}{\partial q^2}\right|_0.
 \label{eq:finite-hencky-reference-biot-tensor}
\end{equation}
Thus directional pressure coupling arises from a mixed mineral-volume--skeleton
strain tangent, without assigning an independent mineral-shape equilibrium.

\section{Generalized-Hooke specialization}
\label{sec:hooke}

For \(\F=\I+\nabla\bm u\), let \(\bm e=\sym\nabla\bm u\).  The
linearization of \eqref{eq:constrained-mineral-hencky-strain} is
\begin{equation}
 \bm\varepsilon=\bm e+O(\lVert\nabla\bm u\rVert^2),
 \qquad
 \bm\eta=\Pdev:\bm e+\frac{q}{3}\I+o(\lVert\nabla\bm u\rVert+|q|).
 \label{eq:hencky-small-strain-limit}
\end{equation}
To expose the mineral-stiffness contribution, begin with the quadratic
Hencky tangent
\begin{equation}
 W_{\mathrm H,0}(\bm e,\bm\eta)=
 \frac12\bm e:\C:\bm e+
 \bm e:\D:\bm\eta+
 \frac12\bm\eta:\Hh:\bm\eta,
 \label{eq:underlying-quadratic-energy}
\end{equation}
Here \(\C\) and \(\Hh\) are fourth-order elasticity tensors: \(\C\) is
the skeleton tensor, whereas \(\Hh\) represents mineral-solid elasticity in
an energy measured per unit reference mixture volume.  The latter
normalization is a volume scaling, not an additional mineral property.  Let
\(\bar W_{m,0}(\bm\eta)\) be the mineral energy per unit reference mineral
volume and let \(\bar{\mathbb H}\) be its fourth-order elasticity tensor.
For a uniform reference
mixture, its contribution to the mixture-reference energy and its tangent are
\begin{equation}
 W_{m,0}=\phi_{s0}\bar W_{m,0},
\qquad
 \Hh=\phi_{s0}\bar{\mathbb H},
\qquad
 \bar{\mathbb H}
 =\left.\frac{\partial^2\bar W_{m,0}}
 {\partial\bm\eta\,\partial\bm\eta}\right|_{\bm\eta=\bm0}.
 \label{eq:mixture-reference-mineral-elasticity}
\end{equation}
Thus \(\bar{\mathbb H}\), rather than \(\Hh\), is the tensor that can
be identified from an isolated mineral specimen, provided that specimen is
represented by the same mineral constitutive state.  The tensor \(\Hh\)
then places that measured response on the reference volume of the porous
mixture.  Both \(\C\) and \(\Hh\) are self-adjoint on symmetric tensors.
The fourth-order tensor \(\D\) maps mineral strain to skeleton stress and
represents their elastic coupling.  The constrained relation in
\eqref{eq:hencky-small-strain-limit} reduces
\eqref{eq:underlying-quadratic-energy} to
\begin{equation}
 W_0(\bm e,q)=
 \frac12\bm e:\C^c:\bm e+
 q\,\bm g:\bm e+
 \frac12 hq^2,
 \label{eq:quadratic-energy}
\end{equation}
where \(\C^c\) is the constrained skeleton tangent, \(\bm g\) is the
mineral-volume--skeleton coupling, and \(h>0\) is the mineral-volume tangent.
Direct substitution gives
\begin{align}
 \C^c&=\C+\sym(\D:\Pdev)+\Pdev:\Hh:\Pdev,
 \label{eq:constrained-skeleton-tangent}\\
 \bm g&=\frac13\left(\D:\I+\Pdev:\Hh:\I\right),
 \label{eq:reduced-volume-strain-coupling}\\
 h&=\frac19\I:\Hh:\I.
 \label{eq:reduced-mineral-volume-tangent}
\end{align}
Here \(\sym(\D:\Pdev)=\tfrac12(\D:\Pdev+\Pdev:\D^T)\).  Thus,
mineral anisotropy enters \(\bm g\) whenever \(\Hh:\I\) has a deviatoric
part.  Mineral equilibrium under pore pressure is
\begin{equation}
 \bm g:\bm e+hq=-\phi_{s0}p,
 \label{eq:linear-mineral-equilibrium}
\end{equation}
and hence
\begin{equation}
 q=-h^{-1}\bm g:\bm e-h^{-1}\phi_{s0}p.
 \label{eq:mineral-strain-solution}
\end{equation}
At \(p=0\), elimination of \(q\) gives the drained skeleton stiffness,
\begin{equation}
 \C^d=\C^c-h^{-1}\bm g\otimes\bm g.
 \label{eq:drained-stiffness}
\end{equation}
Here \((\bm g\otimes\bm g):\bm e=\bm g(\bm g:\bm e)\).
Since \(\bar J=1+q+o(|q|)\), \eqref{eq:mineral-strain-solution}
and \eqref{eq:finite-biot-tensor} yield the small-strain Biot tensor
\begin{equation}
 \boxed{\quad
 \bm B_0=\I+\phi_{s0}h^{-1}\bm g.
 \quad}
 \label{eq:linear-biot-tensor}
\end{equation}
The stress relation is \(\bm\sigma=\bm\sigma^{\prime\prime}-p\bm B_0\).
Equation~\eqref{eq:linear-biot-tensor} shows that \(\C^d\) and \(h\) alone
do not determine directional pressure coupling: \(\bm g\) remains additional
constitutive data.

\section{Special cases}
\label{sec:special-cases}

\subsection{Isotropic mineral stiffness and anisotropic skeleton coupling}

Let the mineral stiffness be isotropic,
\begin{equation}
 \Hh=\phi_{s0}\left[
 \lambda_s\,\I\otimes\I+2\mu_s\,\mathbb I_{\mathrm{sym}}
 \right],
 \qquad K_s=\lambda_s+\frac23\mu_s.
 \label{eq:isotropic-mineral-stiffness}
\end{equation}
Then \(\Pdev:\Hh:\I=\bm0\) and \(h=\phi_{s0}K_s\), so
\begin{equation}
 \bm g=\frac13\D:\I,
 \qquad
 \bm B_0=\I+\frac{1}{3K_s}\D:\I.
 \label{eq:isotropic-mineral-biot-tensor}
\end{equation}
The constrained and drained skeleton tangents may be anisotropic.  With an
isotropic mineral, directional pressure coupling is carried by the volumetric
part \(\D:\I\) of the skeleton--mineral coupling.

\subsection{Anisotropic mineral stiffness with isotropic coupling}

Let the coupling be isotropic,
\begin{equation}
 \D=d_v\Pvol+d_d\Pdev.
 \label{eq:isotropic-skeleton-mineral-coupling}
\end{equation}
For an arbitrary positive-definite mineral stiffness \(\Hh\), equations
\eqref{eq:reduced-volume-strain-coupling}--\eqref{eq:linear-biot-tensor}
give
\begin{equation}
 \bm B_0=
 \I+\frac{\phi_{s0}}{3h}
 \left(d_v\I+\Pdev:\Hh:\I\right),
 \qquad h=\frac19\I:\Hh:\I.
 \label{eq:anisotropic-mineral-biot-tensor}
\end{equation}
Consequently, an anisotropic mineral stiffness can produce a directional Biot
tensor even when the skeleton--mineral coupling is isotropic.  The directional
part is proportional to \(\Pdev:\Hh:\I\), which measures
volumetric--deviatoric coupling in the mineral stiffness.

\subsection{Mineral anisotropy without volumetric--deviatoric coupling}

If the anisotropic mineral stiffness satisfies
\begin{equation}
 \Pdev:\Hh:\I=\bm0,
 \label{eq:mineral-volumetric-deviatoric-decoupling}
\end{equation}
then the mineral may remain anisotropic in its deviatoric block
\(\Pdev:\Hh:\Pdev\), but it does not produce directional pressure coupling
through \eqref{eq:anisotropic-mineral-biot-tensor}.  For the isotropic
coupling \eqref{eq:isotropic-skeleton-mineral-coupling},
\begin{equation}
 \bm B_0=\left(1+\frac{\phi_{s0}d_v}{3h}\right)\I.
 \label{eq:decoupled-mineral-biot-tensor}
\end{equation}
Thus, this class of mineral anisotropy changes the scalar pressure response
through \(h\), but cannot by itself make \(\bm B_0\) nonspherical.

\subsection{Volumetric calibration for the classical coefficient}
\label{sec:volumetric-calibration}

The quadratic family includes a volumetric calibration that recovers the
classical scalar Biot coefficient.  Take the isotropic mineral stiffness
\eqref{eq:isotropic-mineral-stiffness} and set
\begin{equation}
\D=-3K\Pvol+d_d\Pdev,
\qquad
\C^c=\C^d+\frac{K^2}{\phi_{s0}K_s}\I\otimes\I.
 \label{eq:volumetric-energy}
\end{equation}
Here \(K\) and \(K_s\) are the drained skeleton and mineral bulk moduli.
Equations \eqref{eq:reduced-volume-strain-coupling} and
\eqref{eq:reduced-mineral-volume-tangent} give \(\bm g=-K\I\) and
\(h=\phi_{s0}K_s\).  The Schur complement in
\eqref{eq:drained-stiffness} then returns \(\C^d\), and
\eqref{eq:linear-biot-tensor} gives
\begin{equation}
 \bm B_0=\left(1-\frac{K}{K_s}\right)\I,
 \qquad B=1-\frac{K}{K_s}.
 \label{eq:classical-biot}
\end{equation}
The deviatoric coefficient \(d_d\) does not enter \(\bm B_0\), but it
contributes to the constrained tangent in
\eqref{eq:constrained-skeleton-tangent}.
Positive bulk and shear moduli are required for isotropic elastic response.  If
the drained skeleton is no stiffer in bulk than its unjacketed mineral
response, \(0\leq K\leq K_s\), then \eqref{eq:classical-biot} gives
\(0\leq B\leq1\) \citep{makhnenkolabuz2016}.  This modulus ordering is a
material restriction and is not a consequence of tensor anisotropy alone.

\section{Discussion and conclusions}
\label{sec:conclusions}

In our earlier work, we used a pressure Legendre transform to derive a
finite-deformation scalar Biot coefficient \citep{fosterxu2025}.  Here, we
show that fixed pore pressure produces a tensorial Biot correction when true
mineral volume responds directionally to skeleton deformation.  The
distention relation \eqref{eq:distention} remains the kinematic link between
intrinsic density, solid volume fraction, and skeleton volume.  The
isotropic-distention split constrains true-mineral shape through the skeleton
deformation and mineral volume ratio, while directional response is carried by
the mineral-volume equilibrium's constitutive dependence on skeleton strain.

The generalized-Hooke model identifies \(\bm g\) as the additional
constitutive tangent and shows that an anisotropic mineral stiffness enters it
through \(\Pdev:\Hh:\I\).  It recovers
\eqref{eq:classical-biot} under the explicit volumetric calibration
\eqref{eq:volumetric-energy}.  At the small-strain reference state, \(\bm g\)
is observable from the Biot tensor when \(h\) is independently known, since
\eqref{eq:linear-biot-tensor} gives
\(\bm g=\phi_{s0}^{-1}h(\bm B_0-\I)\).  This reduced observable does not
identify the full skeleton--mineral coupling tensor.  Given \(\bm g\) and
\(\Hh\), \eqref{eq:reduced-volume-strain-coupling} fixes only
\(\D:\I=3\bm g-\Pdev:\Hh:\I\); any change to \(\D\) with zero contraction
against \(\I\) leaves \(\bm g\) unchanged.  A
finite-deformation implementation then requires an objective hyperelastic
energy, a locally solvable scalar mineral-volume stationarity problem
\eqref{eq:mineral-coordinate-equilibrium}, and verification of the resulting
tangent and stress transform.  Together, these results provide the
constitutive foundation for an anisotropic numerical model.

\bibliographystyle{plainnat}
\bibliography{references}
\end{document}
